% Laget automatisk av elevark/lag_elevark.py fra sim/samtidighet.html. Endringer her blir overskrevet.
\documentclass[11pt]{article}
\usepackage{elevark}
\begin{document}
\elevarkhode{Fysikk 2 · Spesiell relativitetsteori}{Toget og lynene}{https://fysikklab.no/sim/samtidighet.html}{Simuleringen}

Et 300~m langt tog kjører forbi en perrong. To lyn slår ned i hver sin ende av toget og lager merker både på toget og på perrongen. P står på perrongen midt mellom merkene. T sitter midt i toget. Slo lynene ned samtidig?

\section*{Oppgaver}

Lyset har farten $c$ i begge referansesystemene. Det er det eneste dere trenger for å løse oppgavene.

\begin{oppgaveliste}
\oppgave \fase{gjett}{Gjett}Lynene slår ned samtidig sett fra perrongen. Ser P de to blinkene samtidig? Hva med T, som sitter midt i toget?
\svar
\oppgave \fase{test}{Test}Kjør simuleringen i perrongens system. Hvem ser hvilket blink først?
\svar
\oppgave \fase{forklar}{Forklar}T sitter midt i toget og vet at lyset går like fort i begge retninger. Likevel ser T blinket fra forenden først. Hva må T konkludere om rekkefølgen på lynnedslagene?
\svar
\oppgave \fase{test}{Test}Bytt til togets system og kjør igjen. Stemmer det dere kom fram til? Hva har skjedd med lengden til toget og perrongen?
\svar
\oppgave \fase{forklar}{Regn}Bruk tidene i tabellen for togets system: Når ser T de to blinkene? Lyset går 150~m fra hver ende. Regn ut når lynene slo ned i togets system.
\svar
\oppgave \fase{gjett}{Gjett}Hvilke ting er P og T enige om, og hvilke er de uenige om? Se på hvem som ser hva først.
\svar
\oppgave \fase{utfordring}{Utfordring}Kan to hendelser bytte rekkefølge i et annet referansesystem hvis den ene er årsaken til den andre? Tips: Kunne et signal ha gått fra det ene lynnedslaget til det andre?
\svar
\end{oppgaveliste}

\end{document}
